\documentclass[a4paper]{article}
\usepackage{a4wide,amssymb,epsfig,latexsym,multicol,array,hhline,fancyhdr}
\usepackage{vntex}
\usepackage{amsmath}
\usepackage{lastpage}
\usepackage[lined,boxed,commentsnumbered]{algorithm2e}
\usepackage{enumerate}
\usepackage{color}
\usepackage{graphicx}							% Standard graphics package
\usepackage{array}
\usepackage{tabularx, caption}
\usepackage{multirow}
\usepackage{multicol}
\usepackage{rotating}
\usepackage{graphics}
\usepackage{geometry}
\usepackage{setspace}
\usepackage{epsfig}
\usepackage{tikz}
\usetikzlibrary{arrows,snakes,backgrounds}
\usepackage{hyperref}
\hypersetup{urlcolor=blue,linkcolor=black,citecolor=black,colorlinks=true}
%\usepackage{pstcol} 								% PSTricks with the standard color package

\newtheorem{theorem}{{\bf Theorem}}
\newtheorem{property}{{\bf Property}}
\newtheorem{proposition}{{\bf Proposition}}
\newtheorem{corollary}[proposition]{{\bf Corollary}}
\newtheorem{lemma}[proposition]{{\bf Lemma}}

\AtBeginDocument{\renewcommand*\contentsname{Contents}}
\AtBeginDocument{\renewcommand*\refname{References}}
%\usepackage{fancyhdr}
\setlength{\headheight}{40pt}
\pagestyle{fancy}
\fancyhead{} % clear all header fields
\fancyhead[L]{
 \begin{tabular}{rl}
    \begin{picture}(25,15)(0,0)
    \put(0,-8){\includegraphics[width=8mm, height=8mm]{hcmut.png}}
    %\put(0,-8){\epsfig{width=10mm,figure=hcmut.eps}}
   \end{picture}&
	%\includegraphics[width=8mm, height=8mm]{hcmut.png} & %
	\begin{tabular}{l}
		\textbf{\bf \ttfamily University of Technology, Ho Chi Minh City}\\
		\textbf{\bf \ttfamily Faculty of Computer Science and Engineering}
	\end{tabular}
 \end{tabular}
}
\fancyhead[R]{
	\begin{tabular}{l}
		\tiny \bf \\
		\tiny \bf
	\end{tabular}  }
\fancyfoot{} % clear all footer fields
\fancyfoot[L]{\scriptsize \ttfamily Assignment for Discrete Structures for Computing - Academic year 2019 - 2020}
\fancyfoot[R]{\scriptsize \ttfamily Page {\thepage}/\pageref{LastPage}}
\renewcommand{\headrulewidth}{0.3pt}
\renewcommand{\footrulewidth}{0.3pt}


%%%
\setcounter{secnumdepth}{4}
\setcounter{tocdepth}{3}
\makeatletter
\newcounter {subsubsubsection}[subsubsection]
\renewcommand\thesubsubsubsection{\thesubsubsection .\@alph\c@subsubsubsection}
\newcommand\subsubsubsection{\@startsection{subsubsubsection}{4}{\z@}%
                                     {-3.25ex\@plus -1ex \@minus -.2ex}%
                                     {1.5ex \@plus .2ex}%
                                     {\normalfont\normalsize\bfseries}}
\newcommand*\l@subsubsubsection{\@dottedtocline{3}{10.0em}{4.1em}}
\newcommand*{\subsubsubsectionmark}[1]{}
\makeatother


\begin{document}

\begin{titlepage}
\begin{center}
VIETNAM NATIONAL UNIVERSITY, HO CHI MINH CITY \\
UNIVERSITY OF TECHNOLOGY \\
FACULTY OF COMPUTER SCIENCE AND ENGINEERING
\end{center}

\vspace{1cm}

\begin{figure}[h!]
\begin{center}
\includegraphics[width=3cm]{hcmut.png}
\end{center}
\end{figure}

\vspace{1cm}


\begin{center}
\begin{tabular}{c}
\multicolumn{1}{l}{\textbf{{\Large DISCRETE STRUCTURES FOR COMPUTING (CO1007)}}}\\
~~\\
\hline
\\
\multicolumn{1}{l}{\textbf{{\Large Assignment}}}\\
\\
\textbf{{\Huge Relation - Counting - Probability}}\\
\\
\textbf{{\Huge and Graph}}\\
\\
\hline
\end{tabular}
\end{center}

\vspace{3cm}

\begin{table}[h]
\begin{tabular}{rrl}
\hspace{5 cm} & Advisor: & Fullname\\
& Students: & Fullname of Student 1 - Student 1 ID numbers. \\
& & Fullname of Student 2 - Student 2 ID numbers. \\
\end{tabular}
\end{table}

\begin{center}
{\footnotesize HO CHI MINH CITY, SEPTEMBER 2020}
\end{center}
\end{titlepage}


%\thispagestyle{empty}

\newpage
\tableofcontents
\newpage


%%%%%%%%%%%%%%%%%%%%%%%%%%%%%%%%%
\section{Member list \& Workload}

\begin{center}
\begin{tabular}{|c|c|c|l|c|}
\hline
\textbf{No.} & \textbf{Fullname} & \textbf{Student ID} & \textbf{Problems} & \textbf{Percentage of work}\\
\hline
%%%%%Student 1%%%%%%%%%%
\multirow{3}{*}{1} & \multirow{3}{*}{Lưu Quốc Bình} & \multirow{3}{*}{2033009} & - Exercise 1 & \multirow{3}{*}{30\%}\\
 & &  & Bonus: 1, 2, 3. &\\
 & &  & - Probability: 1, 2, 3. &\\
\hline
%%%%%Student 2%%%%%%%%%%%
\multirow{3}{*}{2} & \multirow{3}{*}{Nguyễn Văn B} & \multirow{3}{*}{19181717} & - Relation \& Counting: 4, 5, 6& \multirow{3}{*}{20\%}\\
 & &  & Bonus: 4, 5, 6. &\\
 & &  & - Graph: 1, 2, 3, Bonus: 1, 2, 3. &\\
\hline
%%%%%Student 1%%%%%%%%%%
\multirow{3}{*}{1} & \multirow{3}{*}{Nguyễn Văn A} & \multirow{3}{*}{19181716} & - Relation \& Counting: 1, 2, 3& \multirow{3}{*}{30\%}\\
& &  & Bonus: 1, 2, 3. &\\
& &  & - Probability: 1, 2, 3. &\\
\hline
%%%%%Student 1%%%%%%%%%%
\multirow{3}{*}{1} & \multirow{3}{*}{Nguyễn Văn A} & \multirow{3}{*}{19181716} & - Relation \& Counting: 1, 2, 3& \multirow{3}{*}{30\%}\\
& &  & Bonus: 1, 2, 3. &\\
& &  & - Probability: 1, 2, 3. &\\
\hline
%%%%%Student 3%%%%%%%%%%%
\end{tabular}
\end{center}

%%%%%%%%%%%%%%%%%%%%%%%%%%%%%%%%%
\section{Relation \& Counting}
	\subsection{Problem 1}
	Write on the report a very detailed introduction to the IVPs Sys. (3) and
	the formulae of its exact solutions for general a, b, c, and d and initial condition R0 and J0. Then
	complete Tab. 2 for all possible cases of eigenvalues of general 2 × 2 matrix A
	\title{How hard would it be to build a spaceship from scrap}

	\subsection{Problem 2}
	...

	\subsection{Bonus exercises}
	...

%%%%%%%%%%%%%%%%%%%%%%%%%%%%%%%%%
\section{Probabilty}
	\subsection{Problem 1}
	...

	\subsection{Problem 2}
	...

	\subsection{Bonus exercises}
	...

%%%%%%%%%%%%%%%%%%%%%%%%%%%%%%%%%


\section{Graph}
	\subsection{Problem 1}
	...

	\subsection{Problem 2}
	...

	\subsection{Bonus exercises}
	...

\begin{thebibliography}{80}


\bibitem{bib1}
...


\bibitem{bib2}
...


\end{thebibliography}
\end{document}

